% ============================================================
%  BoviBot — Rapport de Conception et Réalisation
%  ESP/UCAD — Licence 3 DIC2 — Avril 2026
%  Compiler avec pdflatex (deux passes pour tikz)
%  Sur Overleaf : importer MCD_BoviBot.png et MLD_BoviBot.png
% ============================================================

\documentclass[12pt, a4paper, oneside]{report}

% ── Encodage et langue ───────────────────────────────────────
\usepackage[utf8]{inputenc}
\usepackage[T1]{fontenc}
\usepackage[french]{babel}

% ── Police : Palatino (élégante, professionnelle) ────────────
\usepackage{mathpazo}
\usepackage{microtype}

% ── Tolérance de débordement horizontal ──────────────────────
\emergencystretch=3em
\tolerance=1500
\hbadness=1500

% ── Mise en page ─────────────────────────────────────────────
\usepackage[
  a4paper,
  top=2.5cm,
  bottom=2.5cm,
  left=3cm,
  right=2.5cm,
  headheight=15pt
]{geometry}
\usepackage{setspace}
\onehalfspacing

% ── Couleurs BoviBot ─────────────────────────────────────────
\usepackage[dvipsnames, svgnames, table]{xcolor}
\definecolor{bovibrown}{HTML}{7C4F1E}
\definecolor{bovibrowndark}{HTML}{5A3A15}
\definecolor{bovilight}{HTML}{FDF6EC}
\definecolor{bovibg}{HTML}{F5F0E8}
\definecolor{bovired}{HTML}{DC2626}
\definecolor{boviorange}{HTML}{F59E0B}
\definecolor{boviblue}{HTML}{3B82F6}
\definecolor{bovigreen}{HTML}{16A34A}
\definecolor{bovitext}{HTML}{1D1C17}
\definecolor{bovitextsoft}{HTML}{51453A}
\definecolor{codebg}{HTML}{1E293B}
\definecolor{codetext}{HTML}{E2E8F0}
\definecolor{codekw}{HTML}{93C5FD}
\definecolor{codestr}{HTML}{FCD34D}
\definecolor{codecomment}{HTML}{64748B}

% ── Titres et sections ────────────────────────────────────────
\usepackage{titlesec}

% Style des chapitres
\titleformat{\chapter}[display]
  {\normalfont\huge\bfseries\color{bovibrown}}
  {{\normalsize\color{bovitextsoft}\MakeUppercase{\chaptertitlename}}~\thechapter}
  {10pt}
  {\Huge}
  [\vspace{4pt}{\color{bovibrown}\rule{\textwidth}{1.5pt}}]
\titlespacing*{\chapter}{0pt}{-10pt}{30pt}

% Style des sections
\titleformat{\section}
  {\normalfont\Large\bfseries\color{bovibrowndark}}
  {\thesection}
  {1em}{}
\titlespacing*{\section}{0pt}{20pt}{8pt}

% Style des sous-sections
\titleformat{\subsection}
  {\normalfont\large\bfseries\color{bovibrowndark}}
  {\thesubsection}
  {1em}{}
\titlespacing*{\subsection}{0pt}{14pt}{4pt}

% Style des sous-sous-sections
\titleformat{\subsubsection}
  {\normalfont\normalsize\bfseries\color{bovitextsoft}}
  {\thesubsubsection}
  {1em}{}

% ── En-têtes et pieds de page ─────────────────────────────────
\usepackage{fancyhdr}
\pagestyle{fancy}
\fancyhf{}
\fancyhead[L]{\small\color{bovitextsoft}\slshape\leftmark}
\fancyhead[R]{\small\color{bovibrown}\textbf{BoviBot} \color{bovitextsoft}— ESP/UCAD}
\fancyfoot[C]{\small\color{bovitextsoft}\thepage}
\renewcommand{\headrulewidth}{0.5pt}
\renewcommand{\headrule}{{\color{bovibrown}\hrule width\headwidth height\headrulewidth}}
\renewcommand{\footrulewidth}{0pt}

% Style plain pour les premières pages de chapitre
\fancypagestyle{plain}{
  \fancyhf{}
  \fancyfoot[C]{\small\color{bovitextsoft}\thepage}
  \renewcommand{\headrulewidth}{0pt}
}

% ── Tableaux ─────────────────────────────────────────────────
\usepackage{booktabs}
\usepackage{longtable}
\usepackage{array}
\usepackage{tabularx}
\usepackage{multirow}
\usepackage{colortbl}
\newcolumntype{L}[1]{>{\raggedright\arraybackslash}p{#1}}
\newcolumntype{C}[1]{>{\centering\arraybackslash}p{#1}}

% ── Images ───────────────────────────────────────────────────
\usepackage{graphicx}
\usepackage{float}
\usepackage{caption}
\usepackage{subcaption}
\captionsetup{
  font=small,
  labelfont={bf, color=bovibrown},
  textfont={color=bovitextsoft},
  labelsep=period,
  skip=6pt
}
\captionsetup[table]{position=top}

% ── Boîtes colorées ──────────────────────────────────────────
\usepackage{tcolorbox}
\tcbuselibrary{breakable, skins, listings}

% Boîte info (note bleue)
\newtcolorbox{infobox}[1][]{
  enhanced, breakable,
  colback=boviblue!8,
  colframe=boviblue!60,
  arc=4pt, boxrule=0.6pt,
  left=8pt, right=8pt, top=6pt, bottom=6pt,
  fontupper=\small\color{bovitext},
  #1
}

% Boîte avertissement (orange)
\newtcolorbox{warnbox}[1][]{
  enhanced, breakable,
  colback=boviorange!10,
  colframe=boviorange!70,
  arc=4pt, boxrule=0.6pt,
  left=8pt, right=8pt, top=6pt, bottom=6pt,
  fontupper=\small\color{bovitext},
  #1
}

% Boîte critique (rouge)
\newtcolorbox{critbox}[1][]{
  enhanced, breakable,
  colback=bovired!8,
  colframe=bovired!60,
  arc=4pt, boxrule=0.6pt,
  left=8pt, right=8pt, top=6pt, bottom=6pt,
  fontupper=\small\color{bovitext},
  #1
}

% Boîte succès (verte)
\newtcolorbox{successbox}[1][]{
  enhanced, breakable,
  colback=bovigreen!8,
  colframe=bovigreen!60,
  arc=4pt, boxrule=0.6pt,
  left=8pt, right=8pt, top=6pt, bottom=6pt,
  fontupper=\small\color{bovitext},
  #1
}

% ── Code (listings) ──────────────────────────────────────────
\usepackage{listings}

\lstdefinelanguage{MySQL}{
  morekeywords={CREATE,PROCEDURE,FUNCTION,TRIGGER,EVENT,BEGIN,END,%
    DECLARE,START,TRANSACTION,COMMIT,ROLLBACK,INSERT,UPDATE,DELETE,%
    SELECT,FROM,WHERE,INTO,VALUES,SET,IF,THEN,ELSE,ELSEIF,CALL,RETURN,%
    RETURNS,READS,SQL,DATA,DELIMITER,FOR,EACH,ROW,AFTER,BEFORE,ON,%
    SCHEDULE,EVERY,STARTS,DO,CURRENT_TIMESTAMP,DATE_ADD,DATEDIFF,%
    TIMESTAMPDIFF,CURDATE,SIGNAL,SQLSTATE,RESIGNAL,NOT,EXISTS,NULL,%
    AND,OR,IN,LIMIT,ORDER,BY,DESC,ASC,JOIN,LEFT,RIGHT,INNER,GROUP,%
    HAVING,COUNT,SUM,COALESCE,ROUND,CONCAT,INTERVAL,DAY,MONTH,WEEK,%
    INT,DECIMAL,VARCHAR,TEXT,BOOLEAN,TIMESTAMP,ENUM,AUTO_INCREMENT,%
    PRIMARY,KEY,FOREIGN,REFERENCES,UNIQUE,DEFAULT,GLOBAL,EXIT,HANDLER,%
    SQLEXCEPTION,MESSAGE_TEXT,IS,LIKE,DATE,DATE_SUB,DUMPFILE},
  sensitive=false,
  morestring=[b]',
  morestring=[b]",
  morecomment=[l]{--},
  morecomment=[s]{/*}{*/}
}

\lstdefinestyle{bovistyle}{
  language=MySQL,
  backgroundcolor=\color{codebg},
  basicstyle=\ttfamily\footnotesize\color{codetext},
  keywordstyle=\color{codekw}\bfseries,
  stringstyle=\color{codestr},
  commentstyle=\color{codecomment}\itshape,
  numbers=left,
  numberstyle=\tiny\color{codecomment},
  numbersep=10pt,
  breaklines=true,
  breakatwhitespace=true,
  showstringspaces=false,
  tabsize=4,
  keepspaces=true,
  xleftmargin=18pt,
  framexleftmargin=18pt,
  literate=
    {é}{{\'{e}}}1 {è}{{\`{e}}}1 {ê}{{\^{e}}}1 {ë}{{\"e}}1
    {É}{{\'{E}}}1 {È}{{\`{E}}}1 {Ê}{{\^{E}}}1
    {à}{{\`{a}}}1 {â}{{\^{a}}}1
    {À}{{\`{A}}}1
    {î}{{\^{i}}}1 {ï}{{\"i}}1
    {ô}{{\^{o}}}1
    {ù}{{\`{u}}}1 {û}{{\^{u}}}1
    {ç}{{\c{c}}}1 {Ç}{{\c{C}}}1
    {œ}{{\oe}}1
}

\lstdefinestyle{bashstyle}{
  language=bash,
  backgroundcolor=\color{codebg},
  basicstyle=\ttfamily\footnotesize\color{codetext},
  keywordstyle=\color{bovigreen}\bfseries,
  commentstyle=\color{codecomment}\itshape,
  stringstyle=\color{codestr},
  breaklines=true,
  showstringspaces=false,
  xleftmargin=10pt,
  literate=
    {é}{{\'{e}}}1 {è}{{\`{e}}}1 {à}{{\`{a}}}1 {ç}{{\c{c}}}1
}

% Boîte code SQL avec titre
\newtcblisting{sqlbox}[1]{
  listing only,
  listing options={style=bovistyle},
  breakable, enhanced,
  colback=codebg,
  colframe=bovibrown!70,
  arc=4pt, boxrule=0.8pt,
  title={\small\bfseries\color{white}\ttfamily #1},
  attach boxed title to top left={yshift=-2mm, xshift=5mm},
  boxed title style={colback=bovibrown, arc=2pt, boxrule=0pt},
  top=10pt, bottom=4pt,
  left=2pt, right=2pt
}

% Boîte code Bash
\newtcblisting{bashbox}[1]{
  listing only,
  listing options={style=bashstyle},
  breakable, enhanced,
  colback=codebg,
  colframe=bovigreen!50,
  arc=4pt, boxrule=0.8pt,
  title={\small\bfseries\color{white}\ttfamily #1},
  attach boxed title to top left={yshift=-2mm, xshift=5mm},
  boxed title style={colback=bovibrowndark, arc=2pt, boxrule=0pt},
  top=10pt, bottom=4pt,
  left=2pt, right=2pt
}

% ── Liens hypertexte ─────────────────────────────────────────
\usepackage{hyperref}
\hypersetup{
  colorlinks=true,
  linkcolor=bovibrown,
  urlcolor=boviblue,
  citecolor=bovibrown,
  pdftitle={BoviBot --- Gestion d'Elevage Bovin avec IA et PL/SQL},
  pdfauthor={Sall, Ndoye, Kane, Barro --- ESP/UCAD DIC2},
  pdfsubject={Rapport de Conception et Realisation},
  pdfkeywords={MySQL, PL/SQL, FastAPI, LLM, Text-to-SQL, elevage bovin, Docker},
  bookmarksnumbered=true,
  bookmarksopen=true,
  pdfstartview=FitH
}

% ── Annexes ──────────────────────────────────────────────────
\usepackage[toc, page]{appendix}
\renewcommand{\appendixtocname}{Annexes}
\renewcommand{\appendixpagename}{Annexes}

% ── TikZ (page de garde) ─────────────────────────────────────
\usepackage{tikz}
\usetikzlibrary{calc}

% ── Autres utilitaires ───────────────────────────────────────
\usepackage{enumitem}
\usepackage{pifont}
\usepackage{mdframed}
\usepackage{amsmath}
\usepackage{emptypage}

% ── Commandes personnalisées ─────────────────────────────────

% Lien vers une annexe (inline, avec flèche)
\newcommand{\annexeref}[1]{%
  \tcbox[colback=bovilight, colframe=bovibrown, arc=2pt,
         boxsep=1pt, left=3pt, right=3pt, top=1pt, bottom=1pt,
         on line, boxrule=0.5pt]{%
    \small\color{bovibrown}$\hookrightarrow$~\hyperref[#1]{Annexe~\ref{#1}, p.~\pageref{#1}}%
  }%
}

% Check mark vert
\newcommand{\okmark}{\textcolor{bovigreen}{\ding{51}}}

% Séparateur de section
\newcommand{\secsep}{\vspace{4pt}{\color{bovibrown!30}\hrule}\vspace{8pt}}

% Tableau alternant couleurs
\newcommand{\bovirowcolors}{\rowcolors{2}{bovilight}{white}}

% Noms propres
\newcommand{\bovibot}{\textsc{BoviBot}}

% ── Noms français ────────────────────────────────────────────
\renewcommand{\chaptername}{Chapitre}
\addto\captionsfrench{\renewcommand{\appendixname}{Annexe}}
\addto\captionsfrench{\renewcommand{\listfigurename}{Liste des figures}}
\addto\captionsfrench{\renewcommand{\listtablename}{Liste des tableaux}}
\addto\captionsfrench{\renewcommand{\figurename}{Figure}}
\addto\captionsfrench{\renewcommand{\tablename}{Tableau}}

% ─────────────────────────────────────────────────────────────
\begin{document}

% ════════════════════════════════════════════════════════════
%  PAGE DE GARDE
% ════════════════════════════════════════════════════════════
\begin{titlepage}
  \thispagestyle{empty}
  % Neutralise \onehalfspacing sur cette page uniquement
  \setstretch{1}

  % --- Arrière-plan tikz ---
  \begin{tikzpicture}[remember picture, overlay]
    % Bandeau supérieur brun foncé — 10.5 cm pour couvrir titre + en-tête
    \fill[bovibrowndark]
      (current page.north west)
      rectangle
      ([yshift=-10.5cm] current page.north east);

    % Liseré orange sous le bandeau
    \fill[boviorange]
      ([yshift=-10.5cm] current page.north west)
      rectangle
      ([yshift=-10.9cm] current page.north east);

    % Bande inférieure beige chaud
    \fill[bovibg]
      ([yshift=3.2cm] current page.south west)
      rectangle
      (current page.south east);

    % Double liseré brun/orange
    \fill[bovibrown]
      ([yshift=3.2cm] current page.south west)
      rectangle
      ([yshift=2.95cm] current page.south east);
    \fill[boviorange]
      ([yshift=2.95cm] current page.south west)
      rectangle
      ([yshift=2.82cm] current page.south east);

    % Cercles décoratifs coin supérieur droit
    \fill[bovibrown, opacity=0.18]
      ([xshift=-1.2cm, yshift=-1.2cm] current page.north east)
      circle (3.2cm);
    \fill[bovibrown, opacity=0.10]
      ([xshift=-2.5cm, yshift=-2.5cm] current page.north east)
      circle (5.5cm);

    % Cercle décoratif coin inférieur gauche
    \fill[bovibrown, opacity=0.09]
      ([xshift=0.8cm, yshift=0.8cm] current page.south west)
      circle (2.8cm);
  \end{tikzpicture}

  % ── Établissement (sur fond brun foncé) ─────────────────
  \vspace*{0.8cm}
  \begin{center}
    {\color{white}\large\bfseries École Supérieure Polytechnique (ESP)}\\[0.2em]
    {\color{white!80}\normalsize Université Cheikh Anta Diop — Dakar, Sénégal}\\[0.3em]
    {\color{boviorange!80}\rule{0.35\textwidth}{0.6pt}}\\[0.3em]
    {\color{white!65}\small Licence 3 DIC2 \quad\textbullet\quad IABD $\cdot$ SSI $\cdot$ TR}\\[0.1em]
    {\color{white!55}\small\itshape Bases de Données Avancées — Intégration IA \& PL/SQL}
  \end{center}

  \vspace{0.9cm}

  % ── Titre principal ──────────────────────────────────────
  \begin{center}
    {\color{white}\fontsize{54}{60}\selectfont\bfseries
      Bovi\textcolor{boviorange}{Bot}}\\[0.3em]
    {\color{white!80}\large\itshape Gestion d'Élevage Bovin}\\[0.1em]
    {\color{white!80}\large\itshape Intelligence Artificielle \& PL/SQL}\\[0.5em]
    {\color{boviorange}\rule{0.48\textwidth}{1.4pt}}\\[0.25em]
    {\color{white!55}\normalsize\scshape Rapport de Conception et Réalisation}
  \end{center}

  % ── Zone centrale sur fond blanc ─────────────────────────
  \vspace{1.8cm}

  % Membres du groupe
  \begin{center}
    {\large\bfseries\color{bovibrown} Membres du groupe}\\[0.5em]
    \renewcommand{\arraystretch}{1.25}
    \begin{tabular}{@{\hspace{0.3cm}}L{6.8cm}L{5.8cm}@{}}
      \toprule
      \rowcolor{bovibrown!12}
      \textbf{\color{bovibrowndark}Nom complet} &
      \textbf{\color{bovibrowndark}Parcours} \\
      \midrule
      Mouhamadou Madeniyou Sall &
      Intelligence Artificielle \& Big Data \\
      \rowcolor{bovilight}
      Anna Ndoye &
      Sécurité des Systèmes d'Information \\
      Abdoul Aziz Kane &
      Télécommunications \& Réseaux \\
      \rowcolor{bovilight}
      Fatoumata Barro &
      Télécommunications \& Réseaux \\
      \bottomrule
    \end{tabular}
  \end{center}

  \vspace{1.1cm}

  % Encadrant — boîte centrée
  \begin{center}
    \begin{tcolorbox}[
      enhanced,
      colback=bovilight,
      colframe=bovibrown,
      arc=5pt, boxrule=0.7pt,
      width=0.58\textwidth,
      halign=center,
      top=5pt, bottom=5pt, left=10pt, right=10pt
    ]
      {\small\color{bovitextsoft}\scshape Sous la direction de}\\[0.3em]
      {\color{bovibrown}\large\bfseries Pr. Ahmath Bamba MBACKE}\\[0.15em]
      {\small\color{bovitextsoft}\itshape ahmathbamba.mbacke@esp.sn}
    \end{tcolorbox}
  \end{center}

  % ── Bas de page (bande beige) ────────────────────────────
  \vfill

  \begin{center}
    {\color{bovibrowndark}\normalsize
      Année académique \textbf{2025 – 2026}
      \quad\textcolor{bovibrown}{$\boldsymbol{\mid}$}\quad
      Avril \textbf{2026}}
  \end{center}

  \vspace{0.8cm}
\end{titlepage}

% ════════════════════════════════════════════════════════════
%  PAGES PRÉLIMINAIRES — numérotation romaine
% ════════════════════════════════════════════════════════════
\pagenumbering{roman}
\setcounter{page}{1}

% --- Table des matières ---
\tableofcontents
\newpage

% --- Résumé Général ---
\phantomsection
\addcontentsline{toc}{chapter}{Résumé Général}
\chapter*{Résumé Général}

\begin{tcolorbox}[
  enhanced, breakable,
  colback=bovilight,
  colframe=bovibrown,
  arc=5pt, boxrule=0.8pt,
  left=14pt, right=14pt, top=10pt, bottom=10pt
]
\bovibot{} est une application web de gestion d'élevage bovin développée dans le cadre
du cours d'intégration de l'IA et des bases de données avancées (Licence~3, ESP/UCAD).
Le projet répond à une problématique concrète du contexte sénégalais : permettre à un
éleveur, sans formation informatique avancée, de gérer son troupeau \textbf{en langage
naturel} grâce à un assistant LLM connecté à un moteur PL/SQL avancé.
\end{tcolorbox}

\vspace{1em}

\noindent Le système repose sur \textbf{trois piliers complémentaires} :

\vspace{0.4em}

\begin{enumerate}[leftmargin=1.5cm, itemsep=6pt]
  \item \textbf{Base de données MySQL normalisée} (11 tables, 3NF) — moteur PL/SQL
    complet : 3~procédures stockées transactionnelles, 4~fonctions métier,
    4~triggers de vigilance automatique et 3~events planifiés.

  \item \textbf{Backend Python FastAPI} (29 endpoints REST) — orchestration d'un LLM
    (DeepSeek / gpt-4o-mini) via une stratégie de prompting Few-Shot à 4~piliers,
    avec double protection contre les injections de prompt et SQL.

  \item \textbf{Frontend HTML/CSS/JS} (9 pages : tableau de bord, chat, troupeau,
    santé, généalogie, gestation, rapports, stocks, paramètres) — déployé sur VPS
    Ubuntu~22.04 via Docker Compose.
\end{enumerate}

\vspace{0.8em}

\noindent Les 5~cas fonctionnels du Cahier des Charges ont été validés en production,
ainsi que 6~cas limites de sécurité et de robustesse. L'application \textbf{dépasse les
exigences minimales} sur tous les critères évalués : 14~objets PL/SQL (contre 9~requis),
export PDF intégré et double protection sécurité (bonus).

\vspace{1.4em}

\begin{center}
  \colorbox{bovibrown!10}{%
    \begin{minipage}{0.9\textwidth}
      \vspace{0.4em}
      \centering
      {\small\bfseries\color{bovibrown}Mots-clés\quad}
      {\small\color{bovitextsoft}%
        MySQL $\cdot$ PL/SQL $\cdot$ FastAPI $\cdot$ LLM $\cdot$ Text-to-SQL $\cdot$
        Docker $\cdot$ Élevage bovin $\cdot$ Sénégal}
      \vspace{0.4em}
    \end{minipage}%
  }
\end{center}

\newpage

% ════════════════════════════════════════════════════════════
%  CORPS DU RAPPORT — numérotation arabe
% ════════════════════════════════════════════════════════════
\pagenumbering{arabic}
\setcounter{page}{1}

% ────────────────────────────────────────────────────────────
\chapter{Introduction et contexte métier}
% ────────────────────────────────────────────────────────────

\section{Contexte général}

L'élevage bovin occupe une place centrale dans l'économie rurale sénégalaise. Selon les
données du Ministère de l'Élevage, le cheptel bovin national dépasse les \textbf{3,5~millions
de têtes}, constituant une source essentielle de revenus pour les ménages ruraux, de sécurité
alimentaire et de capital social. Pourtant, la majorité des exploitations sont gérées de
manière informelle, sans outil de suivi structuré, ce qui limite la productivité et
l'accès au crédit agricole.

\vspace{0.5em}

\noindent Dans ce contexte, le numérique représente un levier de modernisation à fort impact.
La combinaison de bases de données relationnelles avancées et d'assistants intelligents permet
d'automatiser le suivi sanitaire, la traçabilité des pesées, la gestion de la reproduction
et les alertes critiques — autant de tâches qui reposent aujourd'hui sur la mémoire et les
registres papier de l'éleveur.

\section{Présentation du projet \bovibot{}}

\bovibot{} est une application web de gestion d'élevage bovin développée dans le cadre du
cours d'intégration de l'IA et des bases de données avancées (Licence~3, ESP/UCAD). Elle
intègre deux composantes techniques complémentaires :

\begin{itemize}[itemsep=4pt]
  \item Un \textbf{assistant LLM} capable d'interroger la base de données en langage
    naturel (mode Text-to-SQL) et d'exécuter des actions métier via des procédures
    stockées, avec confirmation explicite obligatoire.
  \item Un \textbf{moteur PL/SQL avancé} comprenant procédures stockées, fonctions,
    triggers et events MySQL Scheduler automatisant les alertes et les calculs de
    croissance du troupeau.
\end{itemize}

L'application s'adresse à un éleveur gérant un troupeau mixte : elle permet de consulter
l'état du troupeau en langage naturel, d'enregistrer des pesées et des ventes via une
interface de chat, et de recevoir des alertes automatiques sur la santé, la vaccination et
les vêlages à venir.

\section{Problématique}

\begin{infobox}
  Comment concevoir un système de gestion d'élevage bovin qui soit à la fois rigoureux
  sur le plan de la logique métier (intégrité des données, contraintes, automatisations)
  et accessible à un utilisateur non technicien via une interface conversationnelle
  en langage naturel ?
\end{infobox}

\section{Objectifs du projet}

\begin{enumerate}[itemsep=4pt]
  \item Modéliser une base de données MySQL normalisée (3NF) couvrant toutes les
    dimensions de la gestion d'un troupeau bovin.
  \item Implémenter les éléments PL/SQL obligatoires : 2 procédures stockées,
    2 fonctions, 3 triggers, 2 events — objectif surpassé avec 14 objets PL/SQL.
  \item Intégrer un assistant LLM capable de générer du SQL à partir de questions en
    langue naturelle et d'appeler les procédures stockées avec double protection
    contre les injections.
  \item Déployer l'application sur une infrastructure accessible en ligne avec tableau
    de bord, gestion des alertes et interface de chat.
\end{enumerate}

% ────────────────────────────────────────────────────────────
\chapter{Modélisation de la base de données}
% ────────────────────────────────────────────────────────────

La base de données \bovibot{} est construite autour de \textbf{11 tables} couvrant
l'intégralité des dimensions de la gestion d'un troupeau bovin. Elle respecte la
troisième forme normale (3NF) pour garantir l'absence de redondance et la cohérence
des données.

\section{Dictionnaire de données}

\subsection*{Conventions}

\begin{center}
  \begin{tabular}{@{}cl@{}}
    \toprule
    \textbf{Symbole} & \textbf{Signification} \\
    \midrule
    PK  & Clé primaire (\texttt{PRIMARY KEY}) \\
    FK  & Clé étrangère (\texttt{FOREIGN KEY}) \\
    NN  & Non nul (\texttt{NOT NULL}) \\
    UQ  & Unique (\texttt{UNIQUE}) \\
    DEF & Valeur par défaut (\texttt{DEFAULT}) \\
    \bottomrule
  \end{tabular}
\end{center}

\vspace{0.5em}

% --- Table races ---
\begin{table}[H]
\caption{Table \texttt{races} — Référentiel des races bovines}
\label{tab:races}
\small\bovirowcolors
\begin{tabularx}{\textwidth}{@{}cL{3.8cm}L{3cm}cX@{}}
  \toprule
  \rowcolor{bovibrown!15}
  \textbf{\#} & \textbf{Champ} & \textbf{Type SQL} & \textbf{C.} & \textbf{Rôle métier} \\
  \midrule
  1 & \texttt{id}                      & \texttt{INT AUTO\_INC}   & PK    & Identifiant technique unique \\
  2 & \texttt{nom}                     & \texttt{VARCHAR(100)}    & NN    & Nom vernaculaire (ex : Zébu Gobra) \\
  3 & \texttt{origine}                 & \texttt{VARCHAR(100)}    & —     & Pays ou région d'origine \\
  4 & \texttt{poids\_adulte\_moyen\_kg}& \texttt{DECIMAL(6,2)}    & —     & Poids adulte moyen — référence croissance \\
  5 & \texttt{production\_lait\_l\_j}  & \texttt{DECIMAL(6,2)}    & DEF~0 & Production laitière journalière (0 = allaitant) \\
  \bottomrule
\end{tabularx}
\end{table}

% --- Table animaux ---
\begin{table}[H]
\caption{Table \texttt{animaux} — Table centrale du troupeau}
\label{tab:animaux}
\small\bovirowcolors
\begin{tabularx}{\textwidth}{@{}cL{3.4cm}L{4.2cm}cX@{}}
  \toprule
  \rowcolor{bovibrown!15}
  \textbf{\#} & \textbf{Champ} & \textbf{Type SQL} & \textbf{C.} & \textbf{Rôle métier} \\
  \midrule
  1  & \texttt{id}              & \texttt{INT AUTO\_INC}                              & PK      & Identifiant interne unique \\
  2  & \texttt{numero\_tag}     & \texttt{VARCHAR(30)}                                & NN,~UQ  & Boucle auriculaire terrain (ex : TAG-001) \\
  3  & \texttt{nom}             & \texttt{VARCHAR(100)}                               & —       & Nom usuel de l'éleveur \\
  4  & \texttt{race\_id}        & \texttt{INT}                                        & FK      & Race → \texttt{races(id)} \\
  5  & \texttt{sexe}            & \texttt{ENUM('M','F')}                              & NN      & Sexe biologique \\
  6  & \texttt{date\_naissance} & \texttt{DATE}                                       & NN      & Base de calcul pour \texttt{fn\_age\_en\_mois} \\
  7  & \texttt{poids\_actuel}   & \texttt{DECIMAL(6,2)}                               & —       & Mis à jour par \texttt{sp\_enregistrer\_pesee} \\
  8  & \texttt{statut}          & \texttt{ENUM('actif','vendu',}\newline\texttt{'mort','quarantaine')} & DEF & État opérationnel \\
  9  & \texttt{mere\_id}        & \texttt{INT}                                        & FK,NULL & Réflexive → mère (même table) \\
  10 & \texttt{pere\_id}        & \texttt{INT}                                        & FK,NULL & Réflexive → père (même table) \\
  11 & \texttt{notes}           & \texttt{TEXT}                                       & —       & Observations libres \\
  12 & \texttt{created\_at}     & \texttt{TIMESTAMP}                                  & DEF     & Horodatage d'enregistrement \\
  \bottomrule
\end{tabularx}
\end{table}

% --- Tables 3-11 condensées ---
Les tables suivantes (3 à 11) sont documentées selon le même modèle.

\begin{table}[H]
\caption{Tables \texttt{pesees}, \texttt{sante}, \texttt{reproduction} — Suivi zootechnique et sanitaire}
\label{tab:pesees_sante_repro}
\small\bovirowcolors
\begin{tabularx}{\textwidth}{@{}L{2.8cm}L{3cm}L{3cm}X@{}}
  \toprule
  \rowcolor{bovibrown!15}
  \textbf{Table} & \textbf{Clé primaire} & \textbf{Clé étrangère} & \textbf{Champs clés et rôle} \\
  \midrule
  \texttt{pesees}       & \texttt{id}      & \texttt{animal\_id}         & \texttt{poids\_kg}, \texttt{date\_pesee}, \texttt{agent} — Historique des pesées \\
  \texttt{sante}        & \texttt{id}      & \texttt{animal\_id}         & \texttt{type}, \texttt{date\_acte}, \texttt{prochain\_rdv} — Journal vétérinaire \\
  \texttt{reproduction} & \texttt{id}      & \texttt{mere\_id, pere\_id} & \texttt{date\_velage\_prevue}, \texttt{statut} — Gestations \\
  \bottomrule
\end{tabularx}
\end{table}

\begin{table}[H]
\caption{Tables \texttt{alimentation}, \texttt{ventes}, \texttt{alertes} — Coûts, commercial, notifications}
\label{tab:alim_ventes_alertes}
\small\bovirowcolors
\begin{tabularx}{\textwidth}{@{}L{2.8cm}L{3cm}L{3cm}X@{}}
  \toprule
  \rowcolor{bovibrown!15}
  \textbf{Table} & \textbf{Clé primaire} & \textbf{Clé étrangère} & \textbf{Champs clés et rôle} \\
  \midrule
  \texttt{alimentation}    & \texttt{id} & \texttt{animal\_id}      & \texttt{type\_aliment}, \texttt{quantite\_kg}, \texttt{cout\_unitaire\_kg} \\
  \texttt{ventes}          & \texttt{id} & \texttt{animal\_id}      & \texttt{acheteur}, \texttt{prix\_fcfa}, \texttt{date\_vente} \\
  \texttt{alertes}         & \texttt{id} & \texttt{animal\_id} NULL & \texttt{type}, \texttt{niveau}, \texttt{message}, \texttt{traitee} \\
  \bottomrule
\end{tabularx}
\end{table}

\begin{table}[H]
\caption{Tables étendues — Audit, lait, stocks}
\label{tab:extended}
\small\bovirowcolors
\begin{tabularx}{\textwidth}{@{}L{4cm}L{2.8cm}L{2.8cm}X@{}}
  \toprule
  \rowcolor{bovibrown!15}
  \textbf{Table} & \textbf{Clé primaire} & \textbf{Clé étrangère} & \textbf{Rôle} \\
  \midrule
  \texttt{historique\_statut} & \texttt{id} & \texttt{animal\_id} & Journal d'audit ACID des changements de statut \\
  \texttt{production\_lait}   & \texttt{id} & \texttt{animal\_id} & Traites journalières (matin/soir) \\
  \texttt{stocks}             & \texttt{id} & —                   & Inventaire intrants avec seuil d'alerte \\
  \bottomrule
\end{tabularx}
\end{table}

\subsection{Récapitulatif des relations inter-tables}

\begin{table}[H]
\caption{Relations inter-tables de la base de données \bovibot{}}
\label{tab:relations}
\small\bovirowcolors
\begin{tabularx}{\textwidth}{@{}L{4cm}L{2.5cm}L{2.5cm}L{1.2cm}X@{}}
  \toprule
  \rowcolor{bovibrown!15}
  \textbf{Table source} & \textbf{Champ FK} & \textbf{Table cible} & \textbf{Card.} & \textbf{Nature} \\
  \midrule
  \texttt{animaux}           & \texttt{race\_id}   & \texttt{races}   & N,1 & Optionnelle \\
  \texttt{animaux}           & \texttt{mere\_id}   & \texttt{animaux} & N,1 & Réflexive, optionnelle \\
  \texttt{animaux}           & \texttt{pere\_id}   & \texttt{animaux} & N,1 & Réflexive, optionnelle \\
  \texttt{pesees}            & \texttt{animal\_id} & \texttt{animaux} & N,1 & Obligatoire \\
  \texttt{sante}             & \texttt{animal\_id} & \texttt{animaux} & N,1 & Obligatoire \\
  \texttt{reproduction}      & \texttt{mere\_id}   & \texttt{animaux} & N,1 & Obligatoire \\
  \texttt{reproduction}      & \texttt{pere\_id}   & \texttt{animaux} & N,1 & Obligatoire \\
  \texttt{alimentation}      & \texttt{animal\_id} & \texttt{animaux} & N,1 & Obligatoire \\
  \texttt{ventes}            & \texttt{animal\_id} & \texttt{animaux} & N,1 & Obligatoire \\
  \texttt{alertes}           & \texttt{animal\_id} & \texttt{animaux} & N,1 & \textbf{Optionnelle} (NULL — alertes globales) \\
  \texttt{historique\_statut}& \texttt{animal\_id} & \texttt{animaux} & N,1 & Obligatoire \\
  \texttt{production\_lait}  & \texttt{animal\_id} & \texttt{animaux} & N,1 & Obligatoire \\
  \bottomrule
\end{tabularx}
\end{table}

\section{Modèle Conceptuel de Données (MCD)}

Le MCD représente les entités métier et leurs associations avant toute traduction en tables
relationnelles. Il met en évidence les cardinalités et la relation réflexive de généalogie
sur l'entité \textbf{ANIMAL}.

\begin{figure}[H]
  \centering
  \includegraphics[width=0.92\textwidth]{MCD_BoviBot.png}
  \caption{Modèle Conceptuel de Données (MCD) — \bovibot{}}
  \label{fig:mcd}
\end{figure}

\textbf{Entités principales :} RACE (référentiel stable), ANIMAL (entité centrale),
PESEE / SANTE / ALIMENTATION / PRODUCTION\_LAIT (événements liés à un animal),
REPRODUCTION (association ternaire mère-père avec attributs propres), VENTE
(transaction définitive), ALERTE et HISTORIQUE\_STATUT (journalisation PL/SQL), STOCK
(inventaire indépendant).

\section{Modèle Logique de Données (MLD)}

Le MLD traduit le MCD en schéma relationnel. Les clés étrangères matérialisent les
associations et les contraintes d'intégrité référentielle sont définies explicitement.

\begin{figure}[H]
  \centering
  \includegraphics[width=0.92\textwidth]{MLD_BoviBot.png}
  \caption{Modèle Logique de Données (MLD) — \bovibot{}}
  \label{fig:mld}
\end{figure}

\subsection*{Justifications des choix de modélisation non-triviaux}

\paragraph{1. Clés étrangères auto-référencées (\texttt{mere\_id}, \texttt{pere\_id})}
La table \texttt{animaux} contient deux FK pointant vers elle-même, permettant de
naviguer l'arbre généalogique sans table intermédiaire. Les deux champs sont nullable
(parenté inconnue possible).

\paragraph{2. Table \texttt{historique\_statut} pour la traçabilité ACID}
Plutôt que d'écraser le champ \texttt{statut}, une table dédiée archive chaque
transition (ancien → nouveau statut avec horodatage), alimentée \emph{exclusivement}
par le trigger \texttt{trg\_historique\_statut}. En cas de litige commercial,
l'éleveur peut reconstituer l'historique complet de chaque animal.

\paragraph{3. Table \texttt{alertes} polyvalente}
La table centralise les alertes automatiques (triggers, events) \emph{et} les alertes
manuelles. Le champ \texttt{animal\_id} est nullable pour les rapports globaux
(résumé hebdomadaire du troupeau). Le champ \texttt{traitee} (BOOLEAN) permet à
l'interface de filtrer les alertes actives sans jamais supprimer l'historique.

% ────────────────────────────────────────────────────────────
\chapter{Éléments PL/SQL — Description et justification métier}
% ────────────────────────────────────────────────────────────

Le moteur PL/SQL de \bovibot{} repose sur \textbf{3 procédures stockées},
\textbf{4 fonctions}, \textbf{4 triggers} et \textbf{3 events MySQL Scheduler}
— soit 14 objets SQL (contre 9 requis par le CDC). Ces éléments constituent la couche
logique métier de la base de données.

\section{Procédures stockées}

Les procédures encapsulent les opérations d'écriture critiques dans des transactions
ACID (START TRANSACTION / COMMIT / ROLLBACK) et sont les \emph{seuls} points d'entrée
que le LLM peut appeler pour modifier les données.

\subsection{\texttt{sp\_enregistrer\_pesee}}

\textbf{Rôle :} Enregistre une pesée, met à jour \texttt{poids\_actuel}, calcule le GMQ
par rapport à la pesée précédente et insère une alerte si GMQ $<$ 300 g/jour.

\begin{center}\small
  \begin{tabular}{@{}llL{7cm}@{}}
    \toprule
    \textbf{Paramètre} & \textbf{Type} & \textbf{Description} \\
    \midrule
    \texttt{p\_animal\_id} & \texttt{INT}          & Identifiant de l'animal \\
    \texttt{p\_poids\_kg}  & \texttt{DECIMAL(6,2)} & Poids mesuré en kg \\
    \texttt{p\_date}       & \texttt{DATE}         & Date de la pesée \\
    \texttt{p\_agent}      & \texttt{VARCHAR(100)} & Agent ayant effectué la pesée \\
    \bottomrule
  \end{tabular}
\end{center}

\textbf{Justification métier :} L'éleveur ne doit pas calculer manuellement le GMQ
ni mettre à jour le poids courant. Cette procédure garantit la cohérence des données
et automatise la détection des retards de croissance.

\textbf{Comportement ACID :} Un \texttt{EXIT HANDLER FOR SQLEXCEPTION} déclenche un
ROLLBACK complet en cas d'erreur. L'INSERT dans \texttt{pesees} et l'UPDATE de
\texttt{poids\_actuel} sont atomiques.

\textbf{Exemple d'appel :}
\begin{center}
  \texttt{CALL sp\_enregistrer\_pesee(1, 325.00, '2026-04-12', 'BoviBot');}
\end{center}

\medskip\noindent\annexeref{app:sp_enregistrer_pesee}

\subsection{\texttt{sp\_declarer\_vente}}

\textbf{Rôle :} Vérifie que l'animal est actif (sinon SIGNAL SQLSTATE '45000'),
enregistre la vente et change le statut en \texttt{'vendu'}.

\textbf{Justification métier :} Empêche la vente d'animaux déjà vendus, morts ou en
quarantaine. Un animal ne peut pas être vendu deux fois.

\textbf{Exemple d'appel :}
\begin{center}
  \texttt{CALL sp\_declarer\_vente(3, 'Oumar Ba', '771234567', 280000.00, 195.00, '2026-04-12');}
\end{center}

\medskip\noindent\annexeref{app:sp_declarer_vente}

\subsection{\texttt{sp\_rapport\_nutritionnel}}

\textbf{Rôle :} Génère un rapport nutritionnel des 30 derniers jours pour un animal :
quantité totale consommée, coût total, aliment principal, GMQ actuel et coût par kg
de gain. Insère une alerte informative dans \texttt{alertes}.

\textbf{Justification métier :} Ce ratio (coût / kg de prise de poids) oriente les
décisions d'ajustement des rations.

\medskip\noindent\annexeref{app:sp_rapport_nutritionnel}

\section{Fonctions}

Les fonctions encapsulent les calculs métier récurrents. Le SYSTEM\_PROMPT interdit
au LLM de recalculer l'âge, le GMQ ou le coût en dehors des fonctions dédiées.

\subsection{\texttt{fn\_age\_en\_mois}}

\textbf{Retourne :} \texttt{INT} — âge de l'animal en mois entiers.

\textbf{Formule :} $\texttt{TIMESTAMPDIFF(MONTH, date\_naissance, CURDATE())}$

\textbf{Utilisation :} Filtrage veaux $<$ 6 mois, classement par tranche d'âge,
trigger \texttt{trg\_alerte\_poids\_faible}.

\medskip\noindent\annexeref{app:fn_age_en_mois}

\subsection{\texttt{fn\_gmq}}

\textbf{Retourne :} \texttt{DECIMAL(6,3)} — Gain Moyen Quotidien en kg/jour.

\textbf{Formule :}
\[
  \text{GMQ} = \frac{\text{poids\_dernière\_pesée} - \text{poids\_première\_pesée}}{\text{nombre de jours}}
\]

Retourne 0 si une seule pesée ou délai nul.

\medskip\noindent\annexeref{app:fn_gmq}

\subsection{\texttt{fn\_cout\_total\_elevage}}

\textbf{Retourne :} \texttt{DECIMAL(12,2)} — coût total en FCFA.

\textbf{Formule :}
$\sum(\texttt{quantite\_kg} \times \texttt{cout\_unitaire\_kg}) + \sum(\texttt{cout\_sante})$

\medskip\noindent\annexeref{app:fn_cout_total_elevage}

\subsection{\texttt{fn\_rentabilite\_estimee}}

\textbf{Retourne :} \texttt{DECIMAL(12,2)} — rentabilité estimée en FCFA.

\textbf{Formule :}
$(\texttt{poids\_actuel} \times 1300\ \text{FCFA/kg}) - \texttt{fn\_cout\_total\_elevage}()$

Le prix de 1 300 FCFA/kg correspond au cours moyen du bétail sur pied au marché de Dakar.

\medskip\noindent\annexeref{app:fn_rentabilite_estimee}

\section{Triggers}

\begin{table}[H]
\caption{Récapitulatif des triggers \bovibot{}}
\label{tab:triggers}
\small\bovirowcolors
\begin{tabularx}{\textwidth}{@{}L{5.2cm}L{3.2cm}L{2.8cm}X@{}}
  \toprule
  \rowcolor{bovibrown!15}
  \textbf{Trigger} & \textbf{Événement} & \textbf{Condition} & \textbf{Action} \\
  \midrule
  \texttt{trg\_historique\_statut}       & BEFORE UPDATE animaux & statut modifié                   & INSERT → \texttt{historique\_statut} \\
  \texttt{trg\_alerte\_vaccination}      & AFTER INSERT sante    & \texttt{prochain\_rdv} dépassé   & Alerte CRITICAL \\
  \texttt{trg\_alerte\_poids\_faible}    & AFTER INSERT pesees   & veau $<$6 mois, poids $<$60 kg   & Alerte CRITICAL \\
  \texttt{trg\_alerte\_pesee\_manquante} & AFTER INSERT pesees   & sans pesée $>$30 j               & Alerte WARNING \\
  \bottomrule
\end{tabularx}
\end{table}

\subsection{\texttt{trg\_historique\_statut}}
Se déclenche \emph{avant} toute mise à jour sur \texttt{animaux}. Si
\texttt{OLD.statut $\neq$ NEW.statut}, insère une ligne dans \texttt{historique\_statut}.
Ce trigger est l'unique source d'écriture de cette table d'audit.
\medskip\noindent\annexeref{app:trg_historique_statut}

\subsection{\texttt{trg\_alerte\_vaccination}}
Après chaque INSERT dans \texttt{sante}, si \texttt{prochain\_rdv IS NOT NULL AND
prochain\_rdv < CURDATE()}, une alerte CRITICAL est créée. Un rappel déjà dépassé
au moment de la saisie indique un oubli avéré.
\medskip\noindent\annexeref{app:trg_alerte_vaccination}

\subsection{\texttt{trg\_alerte\_poids\_faible}}
Utilise \texttt{fn\_age\_en\_mois()} pour vérifier si l'animal a $\leq$ 6 mois et pèse
$<$ 60 kg, auquel cas une alerte CRITICAL est créée. Un veau sous 60 kg avant 6 mois
présente un risque vital immédiat.
\medskip\noindent\annexeref{app:trg_alerte_poids_faible}

\subsection{\texttt{trg\_alerte\_pesee\_manquante}}
À chaque nouvelle pesée saisie, vérifie l'ensemble du troupeau et signale les animaux
sans pesée depuis plus de 30 jours (anti-doublon journalier intégré).
\medskip\noindent\annexeref{app:trg_alerte_pesee_manquante}

\section{Events MySQL Scheduler}

\begin{table}[H]
\caption{Récapitulatif des events MySQL Scheduler}
\label{tab:events}
\small\bovirowcolors
\begin{tabularx}{\textwidth}{@{}L{5cm}L{2.5cm}X@{}}
  \toprule
  \rowcolor{bovibrown!15}
  \textbf{Event} & \textbf{Fréquence} & \textbf{Logique métier} \\
  \midrule
  \texttt{evt\_alerte\_velages}       & Quotidienne  & Alertes INFO pour les vêlages prévus dans les 7 prochains jours \\
  \texttt{evt\_rapport\_croissance}   & Hebdomadaire & Résumé du nombre d'animaux actifs dans \texttt{alertes} \\
  \texttt{evt\_alerte\_cout\_mensuel} & Mensuelle    & Bilan des coûts d'alimentation du mois écoulé \\
  \bottomrule
\end{tabularx}
\end{table}

Tous les events nécessitent \texttt{SET GLOBAL event\_scheduler = ON} (activé
automatiquement par \texttt{schema.sql}).

\noindent\annexeref{app:events}

% ────────────────────────────────────────────────────────────
\chapter{Architecture technique}
% ────────────────────────────────────────────────────────────

\section{Vue d'ensemble}

\bovibot{} repose sur une architecture trois tiers containerisée, déployée sur un VPS
Ubuntu~22.04. Trois services Docker communiquent sur un réseau interne isolé — la
base de données n'est jamais exposée directement sur l'internet.

\begin{table}[H]
\caption{Couches de l'architecture \bovibot{}}
\label{tab:architecture}
\small\bovirowcolors
\begin{tabularx}{\textwidth}{@{}L{2.5cm}L{3.5cm}X@{}}
  \toprule
  \rowcolor{bovibrown!15}
  \textbf{Couche} & \textbf{Technologie} & \textbf{Rôle} \\
  \midrule
  Présentation & HTML / CSS / JavaScript    & 9 pages : tableau de bord, chat, troupeau, santé, généalogie, gestation, rapports, stocks, paramètres \\
  Proxy        & Nginx (port 8080)          & Reverse proxy, CORS, service des fichiers statiques \\
  Application  & Python FastAPI (port 8002) & 29 routes REST, orchestration LLM, appels procédures stockées, validation SQL \\
  Intelligence & DeepSeek / gpt-4o-mini     & Text-to-SQL, détection d'intention (query / action / info) \\
  Données      & MySQL 8.x + Event Scheduler& Stockage, PL/SQL, triggers, events automatisés \\
  \bottomrule
\end{tabularx}
\end{table}

\section{Réseau Docker interne}

Les services communiquent via un réseau Docker dédié (\texttt{bovibot\_network}).
MySQL n'est jamais exposé sur un port public.

\begin{center}
  \texttt{Internet $\longrightarrow$ Nginx (8080) $\longrightarrow$ backend:8002 $\longrightarrow$ db:3306}
\end{center}

\textbf{Choix techniques Docker :}
\begin{itemize}[itemsep=2pt]
  \item Image backend : \texttt{python:3.11-slim} — empreinte minimale.
  \item \texttt{depends\_on: condition: service\_healthy} — le backend ne démarre que
    lorsque MySQL est prêt.
  \item Volume nommé \texttt{mysql\_data} — les données survivent aux redémarrages.
  \item Variables sensibles (\texttt{DB\_PASSWORD}, \texttt{LLM\_API\_KEY}) via
    fichier \texttt{.env} non versionné.
\end{itemize}

\section{Flux de traitement d'une requête}

\begin{enumerate}[itemsep=4pt]
  \item \textbf{Saisie} de la question dans \texttt{chat.html}.
  \item \textbf{POST /api/chat} → FastAPI : \texttt{sanitize\_input()} bloque
    les patterns d'injection.
  \item \textbf{ask\_llm()} → LLM (SYSTEM\_PROMPT + schéma BDD + 6 derniers messages).
  \item \textbf{LLM retourne} un JSON structuré \texttt{\{type, sql | action, params\}}.
  \item \textbf{Routage} :
    \begin{itemize}
      \item \texttt{type="query"} → \texttt{validate\_sql()} → exécution SELECT MySQL.
      \item \texttt{type="action"} → confirmation affichée → (après "Oui")
        \texttt{CALL sp\_*(...)}.
      \item \texttt{type="info"} → question de clarification à l'utilisateur.
    \end{itemize}
  \item \textbf{Réponse} formatée renvoyée et affichée dans l'interface.
\end{enumerate}

% ────────────────────────────────────────────────────────────
\chapter{Prompt Engineering et intégration LLM}
% ────────────────────────────────────────────────────────────

\section{Stratégie de prompting — Les 4 piliers}

La stratégie retenue est le \textbf{Few-Shot Prompting à 4 blocs} implémentée dans
\texttt{app.py} via la fonction \texttt{get\_system\_prompt()}.

\subsection{Pilier 1 (BLOC 1) — Format de réponse obligatoire}

Le LLM répond \emph{exclusivement} en JSON pur (zéro Markdown), garantissant
l'interopérabilité directe avec le backend FastAPI :

\begin{center}\small
  \begin{tabularx}{\textwidth}{@{}L{2.2cm}X@{}}
    \toprule
    \textbf{Type} & \textbf{Structure JSON} \\
    \midrule
    Consultation & \texttt{\{"type":"query",\ "sql":"SELECT ...",\ "explication":"..."\}} \\
    Action       & \texttt{\{"type":"action",\ "action":"nom\_proc",\ "params":\{...\},\ "confirmation":"..."\}} \\
    Info/Manque  & \texttt{\{"type":"info",\ "sql":null,\ "explication":"..."\}} \\
    \bottomrule
  \end{tabularx}
\end{center}

\subsection{Pilier 2 (BLOC 2) — Règles SQL et abstraction PL/SQL}

Trois règles absolues :
\begin{enumerate}[itemsep=2pt]
  \item Filtrer \texttt{statut='actif'} par défaut.
  \item Utilisation obligatoire des fonctions PL/SQL
    (\texttt{fn\_age\_en\_mois}, \texttt{fn\_gmq}, \texttt{fn\_cout\_total\_elevage},
    \texttt{fn\_rentabilite\_estimee}) — aucun calcul brut en SQL autorisé.
  \item Interdiction de \texttt{DELETE}, \texttt{DROP}, \texttt{INSERT} directs.
    Limite de 100 résultats.
\end{enumerate}

\subsection{Pilier 3 (BLOC 3) — Procédures et mots-clés déclencheurs}

\begin{center}\small
  \begin{tabularx}{\textwidth}{@{}L{5.2cm}X@{}}
    \toprule
    \textbf{Procédure} & \textbf{Mots-clés déclencheurs} \\
    \midrule
    \texttt{sp\_enregistrer\_pesee}    & "enregistre pesée", "pèse", "nouveau poids" \\
    \texttt{sp\_declarer\_vente}       & "déclare vente", "vends", "cède l'animal" \\
    \texttt{sp\_rapport\_nutritionnel} & "rapport nutritionnel", "consommation", "ration" \\
    \bottomrule
  \end{tabularx}
\end{center}

\subsection{Pilier 4 (BLOC 4) — Résolution des identifiants et ambiguïté}

L'éleveur utilise les \texttt{numero\_tag} (ex : TAG-001), pas les \texttt{id} internes.
En action, le LLM passe le tag directement ; le backend résout l'ID.
\textbf{Interdiction des sous-requêtes SQL comme valeur de paramètre.}
Si l'animal est inconnu → type \texttt{info} + question de clarification.

\medskip\noindent Le SYSTEM\_PROMPT complet est disponible en~\annexeref{app:sysprompt}

\section{Exemples de dialogues LLM}

\subsection{Dialogue 1 — Consultation globale}

\textbf{Utilisateur :} \textit{«~Liste tous les animaux actifs avec leur âge et GMQ~»}

\textbf{Type détecté :} \texttt{query}

\textbf{SQL généré :}
\begin{center}\small
\texttt{SELECT a.numero\_tag, a.nom, r.nom AS race, fn\_age\_en\_mois(a.id) AS age\_mois,\\
a.poids\_actuel, fn\_gmq(a.id) AS gmq\_kg\_jour\\
FROM animaux a LEFT JOIN races r ON a.race\_id = r.id\\
WHERE a.statut = 'actif' LIMIT 100}
\end{center}

\textbf{Résultat :} 7 animaux avec âge et GMQ pré-calculés par les fonctions MySQL.

\subsection{Dialogue 2 — Action avec confirmation (pesée)}

\textbf{Utilisateur :} \textit{«~Enregistre une pesée de 325 kg pour TAG-001 aujourd'hui~»}

\begin{enumerate}[itemsep=2pt]
  \item LLM retourne \texttt{type:"action"} + message de confirmation.
  \item Système affiche : \textit{«~Confirmez-vous 325 kg pour TAG-001 le 2026-04-12 ?~»}
  \item Utilisateur répond «~Oui~».
  \item Backend résout TAG-001 → \texttt{animal\_id = 1}, exécute :
    \texttt{CALL sp\_enregistrer\_pesee(1, 325.00, '2026-04-12', 'BoviBot');}
  \item Réponse : \textit{«~Pesée enregistrée. GMQ actuel : 0.357 kg/jour. Aucune alerte.~»}
\end{enumerate}

\subsection{Dialogue 3 — Gestion de l'ambiguïté (vente avec donnée manquante)}

\textbf{Utilisateur :} \textit{«~Déclare la vente de TAG-003 à Oumar Ba pour 280 000 FCFA~»}

Le LLM détecte qu'il manque le poids de vente (requis par \texttt{sp\_declarer\_vente}).
Il retourne \texttt{type:"info"} et pose la question. Après réponse de l'utilisateur,
il relance avec les paramètres complets.

\section{Sécurité — Double protection}

\begin{table}[H]
\caption{Double couche de sécurité LLM — \bovibot{}}
\label{tab:securite}
\small\bovirowcolors
\begin{tabularx}{\textwidth}{@{}L{4cm}L{4cm}X@{}}
  \toprule
  \rowcolor{bovibrown!15}
  \textbf{Couche} & \textbf{Fonction} & \textbf{Mécanisme} \\
  \midrule
  Couche 1 — Entrée & \texttt{sanitize\_input()} & 26+ regex de détection de prompt injection (jailbreak, DAN, bypass…) → HTTP 400 \\
  Couche 2 — Sortie & \texttt{validate\_sql()}   & Le SQL généré par le LLM doit commencer par SELECT. Mots-clés interdits : INSERT, UPDATE, DELETE, DROP, TRUNCATE, ALTER, CREATE, UNION, OUTFILE… → HTTP 400 \\
  \bottomrule
\end{tabularx}
\end{table}

% ────────────────────────────────────────────────────────────
\chapter{Tests}
% ────────────────────────────────────────────────────────────

\section{Tests cas normaux — 5 cas du Cahier des Charges}

\begin{table}[H]
\caption{Résultats des 5 cas CDC validés en production}
\label{tab:tests_normaux}
\small\bovirowcolors
\begin{tabularx}{\textwidth}{@{}cL{3.5cm}L{3.2cm}XC{0.7cm}@{}}
  \toprule
  \rowcolor{bovibrown!15}
  \textbf{\#} & \textbf{Question utilisateur} & \textbf{Type détecté} & \textbf{Résultat} & \\
  \midrule
  CDC-01 & Liste animaux actifs + âge + GMQ   & \texttt{query}                          & 7 animaux, \texttt{fn\_gmq()} utilisé              & \okmark \\
  CDC-02 & Animaux GMQ $<$ 0.3 kg/j           & \texttt{query}                          & TAG-006 (0.25), TAG-007 (0.18)                     & \okmark \\
  CDC-03 & Femelles vêlant dans 30 j           & \texttt{query}                          & TAG-002 — vêlage 2026-05-10                        & \okmark \\
  CDC-04 & Pesée 325 kg pour TAG-001           & \texttt{action} + confirmation          & Pesée insérée, GMQ calculé                         & \okmark \\
  CDC-05 & Vente TAG-003, 280\,000 F           & \texttt{info} → \texttt{action}         & Vente enregistrée, statut vendu                    & \okmark \\
  \bottomrule
\end{tabularx}
\end{table}

\section{Tests cas limites}

\begin{table}[H]
\caption{Résultats des 6 cas limites de sécurité et de robustesse}
\label{tab:tests_limites}
\small\bovirowcolors
\begin{tabularx}{\textwidth}{@{}L{1.2cm}L{3.2cm}L{3.5cm}XC{0.7cm}@{}}
  \toprule
  \rowcolor{bovibrown!15}
  \textbf{Cas} & \textbf{Scénario} & \textbf{Comportement attendu} & \textbf{Mécanisme} & \\
  \midrule
  CL-01 & Veau 4 mois, 45 kg             & Alerte CRITICAL automatique               & \texttt{trg\_alerte\_poids\_faible}  & \okmark \\
  CL-02 & Vente animal déjà vendu         & SQLSTATE '45000' + ROLLBACK               & \texttt{sp\_declarer\_vente} guard   & \okmark \\
  CL-03 & RDV vaccination dépassé         & Alerte CRITICAL créée                     & \texttt{trg\_alerte\_vaccination}    & \okmark \\
  CL-04 & Prompt injection (jailbreak)    & HTTP 400 avant transmission au LLM        & \texttt{sanitize\_input()}           & \okmark \\
  CL-05 & Injection SQL via chat          & HTTP 400 — mot-clé INSERT interdit        & \texttt{validate\_sql()}             & \okmark \\
  CL-06 & Animal inexistant (TAG-999)     & Type \texttt{info} + demande clarification& BLOC~4 du SYSTEM\_PROMPT            & \okmark \\
  \bottomrule
\end{tabularx}
\end{table}

% ────────────────────────────────────────────────────────────
\chapter{Guide d'installation et de déploiement}
% ────────────────────────────────────────────────────────────

\section{Prérequis}

\begin{table}[H]
\caption{Prérequis de déploiement \bovibot{}}
\label{tab:prereqs}
\small\bovirowcolors
\begin{tabularx}{\textwidth}{@{}L{4cm}X@{}}
  \toprule
  \rowcolor{bovibrown!15}
  \textbf{Élément} & \textbf{Requis} \\
  \midrule
  OS              & Ubuntu 22.04 LTS (VPS ou machine locale) \\
  Docker + Compose& v24+ (\texttt{curl -fsSL https://get.docker.com | sh}) \\
  Clé API LLM     & OpenAI, DeepSeek, ou tout provider compatible OpenAI SDK \\
  RAM             & 2 Go minimum \\
  Ports ouverts   & 8080 (HTTP), 22 (SSH) \\
  \bottomrule
\end{tabularx}
\end{table}

\section{Installation en 4 commandes}

\begin{bashbox}{Terminal — Déploiement VPS}
# 1. Cloner le dépôt
git clone https://github.com/MADENIYOU/bovibot.git && cd bovibot

# 2. Configurer les variables d'environnement
cp .env.example .env
nano .env          # Renseigner DB_PASSWORD et LLM_API_KEY

# 3. Lancer les 3 services (MySQL + FastAPI + Nginx)
docker compose up -d --build

# 4. Vérifier le déploiement
docker compose ps
\end{bashbox}

\begin{warnbox}
  \textbf{Point critique :} \texttt{DB\_HOST=db} (nom du service Docker dans le réseau
  interne) et \textbf{non} \texttt{localhost}. C'est l'erreur la plus fréquente en
  migration local $\to$ Docker.
\end{warnbox}

\section{Vérifications post-déploiement}

\begin{bashbox}{Terminal — Vérifications}
curl http://localhost:8080/health      # {"status":"ok"}
curl http://localhost:8080/api/animaux # JSON du troupeau
curl http://localhost:8080/api/alertes # Alertes actives
docker compose logs -f backend         # Logs FastAPI temps réel
\end{bashbox}

\section{Activation de l'Event Scheduler MySQL}

Activé automatiquement par \texttt{schema.sql}. Pour vérifier :

\begin{bashbox}{MySQL — Vérification}
docker compose exec db mysql -u root -p bovibot
SHOW VARIABLES LIKE 'event_scheduler';  -- Value: ON
\end{bashbox}

\section{Mise à jour}

\begin{bashbox}{Terminal — Mise à jour}
git pull origin main
docker compose up -d --build
# Volume mysql_data preserve toutes les donnees
\end{bashbox}

% ────────────────────────────────────────────────────────────
\chapter{Conclusion et perspectives}
% ────────────────────────────────────────────────────────────

\section{Bilan des objectifs atteints}

\begin{table}[H]
\caption{Bilan des objectifs par rapport au Cahier des Charges}
\label{tab:bilan}
\small\bovirowcolors
\begin{tabularx}{\textwidth}{@{}L{4.5cm}cX@{}}
  \toprule
  \rowcolor{bovibrown!15}
  \textbf{Critère CDC} & \textbf{} & \textbf{Détail} \\
  \midrule
  8 tables MySQL (3NF)           & \okmark & 11 tables (+ \texttt{historique\_statut}, \texttt{production\_lait}, \texttt{stocks}) \\
  2 procédures stockées          & \okmark & 3 implémentées (+ \texttt{sp\_rapport\_nutritionnel}) \\
  2 fonctions PL/SQL             & \okmark & 4 implémentées (+ \texttt{fn\_cout\_total\_elevage}, \texttt{fn\_rentabilite\_estimee}) \\
  3 triggers                     & \okmark & 4 implémentés (+ \texttt{trg\_alerte\_pesee\_manquante}) \\
  2 events MySQL                 & \okmark & 3 implémentés (+ \texttt{evt\_alerte\_cout\_mensuel}) \\
  LLM Text-to-SQL                & \okmark & 5 cas CDC validés, fonctions PL/SQL intégrées \\
  LLM actions avec confirmation  & \okmark & Flux à 2 étapes, confirmation explicite \\
  Interface web                  & \okmark & 9 pages HTML \\
  Déploiement en ligne           & \okmark & VPS Ubuntu 22.04 via Docker Compose \\
  \textit{Sécurité (bonus)}      & \okmark & Double protection : \texttt{sanitize\_input()} + \texttt{validate\_sql()} \\
  \textit{Export PDF (bonus)}    & \okmark & ReportLab — fiche individuelle + rapport complet \\
  \bottomrule
\end{tabularx}
\end{table}

\section{Limites actuelles}

\begin{itemize}[itemsep=4pt]
  \item \textbf{Pas de notifications push :} les alertes critiques ne déclenchent pas
    d'envoi SMS ou email.
  \item \textbf{Pas d'authentification utilisateur :} toute personne ayant l'URL
    peut accéder aux données.
  \item \textbf{Performance sur grands troupeaux :} les fonctions \texttt{fn\_gmq()} et
    \texttt{fn\_cout\_total\_elevage()} sont appelées ligne par ligne — des index
    supplémentaires seraient nécessaires au-delà de 50 animaux.
  \item \textbf{LLM déterministe :} \texttt{temperature=0} maximise la cohérence mais
    limite la variété des reformulations.
\end{itemize}

\section{Perspectives d'amélioration}

\begin{enumerate}[itemsep=4pt]
  \item \textbf{WebSocket} — remplacer le polling (30 s) par une connexion persistante
    pour les alertes critiques en temps réel.
  \item \textbf{Authentification JWT} — système de login sécurisé.
  \item \textbf{Notifications multi-canaux} — SMS via Orange API Sénégal ou Twilio.
  \item \textbf{Application mobile} — React Native ou Flutter pour un accès terrain.
  \item \textbf{Export Excel} — fichier \texttt{.xlsx} via \texttt{openpyxl}.
  \item \textbf{Analyse prédictive} — prédire la date optimale de vente par animal
    via régression sur le GMQ historique.
  \item \textbf{Optimisation PL/SQL} — colonne calculée pour le GMQ afin d'éviter
    les recalculs répétés sur les grands troupeaux.
\end{enumerate}

% ════════════════════════════════════════════════════════════
%  ANNEXES
% ════════════════════════════════════════════════════════════
\newpage
\appendix
\addappheadtotoc

% ────────────────────────────────────────────────────────────
\chapter{Procédures stockées — Code SQL complet}
\label{app:procedures}
% ────────────────────────────────────────────────────────────

\section{\texttt{sp\_enregistrer\_pesee}}
\label{app:sp_enregistrer_pesee}

\begin{sqlbox}{sp\_enregistrer\_pesee}
DELIMITER $$
CREATE PROCEDURE sp_enregistrer_pesee(
    IN p_animal_id INT,
    IN p_poids_kg  DECIMAL(6,2),
    IN p_date      DATE,
    IN p_agent     VARCHAR(100)
)
BEGIN
    DECLARE v_derniere_pesee DECIMAL(6,2);
    DECLARE v_jours INT;
    DECLARE v_gmq   DECIMAL(6,2);

    -- Gestionnaire ACID : annule tout si erreur SQL
    DECLARE EXIT HANDLER FOR SQLEXCEPTION
    BEGIN
        ROLLBACK;
        RESIGNAL;
    END;

    START TRANSACTION;

        -- 1. Inserer la pesee dans l'historique
        INSERT INTO pesees (animal_id, poids_kg, date_pesee, agent)
        VALUES (p_animal_id, p_poids_kg, p_date, p_agent);

        -- 2. Mettre a jour le poids courant de l'animal
        UPDATE animaux SET poids_actuel = p_poids_kg WHERE id = p_animal_id;

        -- 3. Recuperer la pesee precedente pour calculer le GMQ
        SELECT poids_kg, DATEDIFF(p_date, date_pesee)
        INTO v_derniere_pesee, v_jours
        FROM pesees
        WHERE animal_id = p_animal_id
          AND date_pesee < p_date
        ORDER BY date_pesee DESC
        LIMIT 1;

        -- 4. Alerte si GMQ < 300 g/jour
        IF v_derniere_pesee IS NOT NULL AND v_jours > 0 THEN
            SET v_gmq = (p_poids_kg - v_derniere_pesee) / v_jours;
            IF v_gmq < 0.3 THEN
                INSERT INTO alertes (animal_id, type, message, niveau)
                VALUES (p_animal_id, 'poids',
                    CONCAT('GMQ faible : ', ROUND(v_gmq * 1000),
                           ' g/jour (seuil : 300 g/jour)'),
                    'warning');
            END IF;
        END IF;

    COMMIT;
END$$
DELIMITER ;
\end{sqlbox}

\section{\texttt{sp\_declarer\_vente}}
\label{app:sp_declarer_vente}

\begin{sqlbox}{sp\_declarer\_vente}
DELIMITER $$
CREATE PROCEDURE sp_declarer_vente(
    IN p_animal_id   INT,
    IN p_acheteur    VARCHAR(150),
    IN p_telephone   VARCHAR(20),
    IN p_prix        DECIMAL(12,2),
    IN p_poids_vente DECIMAL(6,2),
    IN p_date_vente  DATE
)
BEGIN
    DECLARE v_statut VARCHAR(20);

    DECLARE EXIT HANDLER FOR SQLEXCEPTION
    BEGIN
        ROLLBACK;
        RESIGNAL;
    END;

    START TRANSACTION;

        -- 1. Verifier que l'animal est en statut 'actif'
        SELECT statut INTO v_statut
        FROM animaux WHERE id = p_animal_id;

        IF v_statut != 'actif' THEN
            SIGNAL SQLSTATE '45000'
                SET MESSAGE_TEXT =
                    'Cet animal ne peut pas etre vendu (statut non actif)';
        END IF;

        -- 2. Enregistrer la vente
        INSERT INTO ventes (animal_id, acheteur, telephone_acheteur,
                            date_vente, poids_vente_kg, prix_fcfa)
        VALUES (p_animal_id, p_acheteur, p_telephone,
                p_date_vente, p_poids_vente, p_prix);

        -- 3. Changer le statut (declenche trg_historique_statut)
        UPDATE animaux SET statut = 'vendu' WHERE id = p_animal_id;

    COMMIT;
END$$
DELIMITER ;
\end{sqlbox}

\section{\texttt{sp\_rapport\_nutritionnel}}
\label{app:sp_rapport_nutritionnel}

\begin{sqlbox}{sp\_rapport\_nutritionnel}
DELIMITER $$
CREATE PROCEDURE sp_rapport_nutritionnel(IN p_animal_id INT)
BEGIN
    DECLARE v_quantite_totale   DECIMAL(10,2);
    DECLARE v_cout_total        DECIMAL(12,2);
    DECLARE v_aliment_principal VARCHAR(100);
    DECLARE v_gmq               DECIMAL(6,3);
    DECLARE v_cout_par_kg_gain  DECIMAL(12,2) DEFAULT 0;

    -- Consommation et cout sur les 30 derniers jours
    SELECT SUM(quantite_kg), SUM(quantite_kg * cout_unitaire_kg)
    INTO v_quantite_totale, v_cout_total
    FROM alimentation
    WHERE animal_id = p_animal_id
      AND date_alimentation >= DATE_SUB(CURDATE(), INTERVAL 30 DAY);

    -- Aliment le plus distribue (en poids)
    SELECT type_aliment INTO v_aliment_principal
    FROM alimentation
    WHERE animal_id = p_animal_id
      AND date_alimentation >= DATE_SUB(CURDATE(), INTERVAL 30 DAY)
    GROUP BY type_aliment
    ORDER BY SUM(quantite_kg) DESC
    LIMIT 1;

    -- Calcul du cout par kg de gain
    SET v_gmq = fn_gmq(p_animal_id);
    IF v_gmq > 0 THEN
        SET v_cout_par_kg_gain = (v_cout_total / 30) / v_gmq;
    END IF;

    -- Alerte informative
    IF v_quantite_totale IS NOT NULL THEN
        INSERT INTO alertes (animal_id, type, message, niveau)
        VALUES (p_animal_id, 'alimentation',
            CONCAT('Rapport 30j: ', ROUND(v_quantite_totale,1), 'kg (',
                   v_aliment_principal, '). Cout: ',
                   ROUND(v_cout_total,0), ' FCFA.'),
            'info');
    END IF;

    -- Retourner les indicateurs
    SELECT
        v_quantite_totale   AS quantite_totale_30j,
        v_cout_total        AS cout_total_30j,
        v_aliment_principal AS aliment_principal,
        v_gmq               AS gmq_actuel,
        v_cout_par_kg_gain  AS cout_kg_gain;
END$$
DELIMITER ;
\end{sqlbox}

% ────────────────────────────────────────────────────────────
\chapter{Fonctions PL/SQL — Code SQL complet}
\label{app:fonctions}
% ────────────────────────────────────────────────────────────

\section{\texttt{fn\_age\_en\_mois}}
\label{app:fn_age_en_mois}

\begin{sqlbox}{fn\_age\_en\_mois}
DELIMITER $$
CREATE FUNCTION fn_age_en_mois(p_animal_id INT)
RETURNS INT
READS SQL DATA
BEGIN
    DECLARE v_date_naissance DATE;
    SELECT date_naissance INTO v_date_naissance
    FROM animaux WHERE id = p_animal_id;
    RETURN TIMESTAMPDIFF(MONTH, v_date_naissance, CURDATE());
END$$
DELIMITER ;
\end{sqlbox}

\section{\texttt{fn\_gmq}}
\label{app:fn_gmq}

\begin{sqlbox}{fn\_gmq}
DELIMITER $$
CREATE FUNCTION fn_gmq(p_animal_id INT)
RETURNS DECIMAL(6,3)
READS SQL DATA
BEGIN
    DECLARE v_premiere_pesee DECIMAL(6,2);
    DECLARE v_derniere_pesee DECIMAL(6,2);
    DECLARE v_premiere_date  DATE;
    DECLARE v_derniere_date  DATE;
    DECLARE v_jours INT;

    SELECT poids_kg, date_pesee INTO v_premiere_pesee, v_premiere_date
    FROM pesees WHERE animal_id = p_animal_id
    ORDER BY date_pesee ASC LIMIT 1;

    SELECT poids_kg, date_pesee INTO v_derniere_pesee, v_derniere_date
    FROM pesees WHERE animal_id = p_animal_id
    ORDER BY date_pesee DESC LIMIT 1;

    SET v_jours = DATEDIFF(v_derniere_date, v_premiere_date);

    -- Retourner 0 si une seule pesee ou aucune
    IF v_jours = 0 OR v_premiere_pesee IS NULL THEN
        RETURN 0;
    END IF;

    RETURN (v_derniere_pesee - v_premiere_pesee) / v_jours;
END$$
DELIMITER ;
\end{sqlbox}

\section{\texttt{fn\_cout\_total\_elevage}}
\label{app:fn_cout_total_elevage}

\begin{sqlbox}{fn\_cout\_total\_elevage}
DELIMITER $$
CREATE FUNCTION fn_cout_total_elevage(p_animal_id INT)
RETURNS DECIMAL(12,2)
READS SQL DATA
BEGIN
    DECLARE v_alim  DECIMAL(12,2) DEFAULT 0;
    DECLARE v_sante DECIMAL(12,2) DEFAULT 0;

    SELECT COALESCE(SUM(quantite_kg * cout_unitaire_kg), 0)
    INTO v_alim FROM alimentation WHERE animal_id = p_animal_id;

    SELECT COALESCE(SUM(cout), 0)
    INTO v_sante FROM sante WHERE animal_id = p_animal_id;

    RETURN v_alim + v_sante;
END$$
DELIMITER ;
\end{sqlbox}

\section{\texttt{fn\_rentabilite\_estimee}}
\label{app:fn_rentabilite_estimee}

\begin{sqlbox}{fn\_rentabilite\_estimee}
DELIMITER $$
CREATE FUNCTION fn_rentabilite_estimee(p_animal_id INT)
RETURNS DECIMAL(12,2)
READS SQL DATA
BEGIN
    DECLARE v_poids          DECIMAL(6,2);
    DECLARE v_cout_total     DECIMAL(12,2);
    DECLARE v_prix_marche_kg DECIMAL(6,2) DEFAULT 1300.00;

    SELECT poids_actuel INTO v_poids FROM animaux WHERE id = p_animal_id;
    SET v_cout_total = fn_cout_total_elevage(p_animal_id);

    IF v_poids IS NULL OR v_cout_total IS NULL THEN
        RETURN 0;
    END IF;

    RETURN (v_poids * v_prix_marche_kg) - v_cout_total;
END$$
DELIMITER ;
\end{sqlbox}

% ────────────────────────────────────────────────────────────
\chapter{Triggers — Code SQL complet}
\label{app:triggers}
% ────────────────────────────────────────────────────────────

\section{\texttt{trg\_historique\_statut}}
\label{app:trg_historique_statut}

\begin{sqlbox}{trg\_historique\_statut}
DELIMITER $$
CREATE TRIGGER trg_historique_statut
BEFORE UPDATE ON animaux
FOR EACH ROW
BEGIN
    IF OLD.statut != NEW.statut THEN
        INSERT INTO historique_statut (animal_id, ancien_statut, nouveau_statut)
        VALUES (OLD.id, OLD.statut, NEW.statut);
    END IF;
END$$
DELIMITER ;
\end{sqlbox}

\section{\texttt{trg\_alerte\_vaccination}}
\label{app:trg_alerte_vaccination}

\begin{sqlbox}{trg\_alerte\_vaccination}
DELIMITER $$
CREATE TRIGGER trg_alerte_vaccination
AFTER INSERT ON sante
FOR EACH ROW
BEGIN
    IF NEW.prochain_rdv IS NOT NULL AND NEW.prochain_rdv < CURDATE() THEN
        INSERT INTO alertes (animal_id, type, message, niveau)
        VALUES (NEW.animal_id, 'vaccination',
            CONCAT('Rappel vaccination en retard depuis le ',
                   NEW.prochain_rdv),
            'critical');
    END IF;
END$$
DELIMITER ;
\end{sqlbox}

\section{\texttt{trg\_alerte\_poids\_faible}}
\label{app:trg_alerte_poids_faible}

\begin{sqlbox}{trg\_alerte\_poids\_faible}
DELIMITER $$
CREATE TRIGGER trg_alerte_poids_faible
AFTER INSERT ON pesees
FOR EACH ROW
BEGIN
    DECLARE v_age_mois INT;
    SELECT fn_age_en_mois(NEW.animal_id) INTO v_age_mois;
    -- Veau de moins de 6 mois : poids < 60 kg = alerte critique
    IF v_age_mois <= 6 AND NEW.poids_kg < 60 THEN
        INSERT INTO alertes (animal_id, type, message, niveau)
        VALUES (NEW.animal_id, 'poids',
            CONCAT('Poids critique pour un veau de ', v_age_mois,
                   ' mois : ', NEW.poids_kg, ' kg'),
            'critical');
    END IF;
END$$
DELIMITER ;
\end{sqlbox}

\section{\texttt{trg\_alerte\_pesee\_manquante}}
\label{app:trg_alerte_pesee_manquante}

\begin{sqlbox}{trg\_alerte\_pesee\_manquante}
DELIMITER $$
CREATE TRIGGER trg_alerte_pesee_manquante
AFTER INSERT ON pesees
FOR EACH ROW
BEGIN
    INSERT INTO alertes (animal_id, type, message, niveau)
    SELECT a.id, 'poids',
           CONCAT('Pesee manquante depuis > 30 jours : ', a.numero_tag),
           'warning'
    FROM animaux a
    WHERE a.statut = 'actif'
      AND NOT EXISTS (
          SELECT 1 FROM pesees p
          WHERE p.animal_id = a.id
            AND p.date_pesee >= DATE_SUB(CURDATE(), INTERVAL 30 DAY)
      )
      AND NOT EXISTS (
          -- Anti-doublon : pas deux alertes du meme type le meme jour
          SELECT 1 FROM alertes al
          WHERE al.animal_id = a.id
            AND al.type = 'poids'
            AND al.message LIKE 'Pesee manquante%'
            AND DATE(al.date_creation) = CURDATE()
      );
END$$
DELIMITER ;
\end{sqlbox}

% ────────────────────────────────────────────────────────────
\chapter{Events MySQL Scheduler — Code SQL complet}
\label{app:events}
% ────────────────────────────────────────────────────────────

\begin{sqlbox}{evt\_alerte\_velages (quotidien)}
DELIMITER $$
CREATE EVENT evt_alerte_velages
ON SCHEDULE EVERY 1 DAY
STARTS CURRENT_TIMESTAMP
DO
BEGIN
    INSERT INTO alertes (animal_id, type, message, niveau)
    SELECT r.mere_id, 'velage',
        CONCAT('Velage prevu dans ',
               DATEDIFF(r.date_velage_prevue, CURDATE()),
               ' jour(s) : ', a.numero_tag),
        'info'
    FROM reproduction r
    JOIN animaux a ON r.mere_id = a.id
    WHERE r.statut = 'en_gestation'
      AND r.date_velage_prevue
          BETWEEN CURDATE() AND DATE_ADD(CURDATE(), INTERVAL 7 DAY)
      AND NOT EXISTS (
          SELECT 1 FROM alertes al
          WHERE al.animal_id = r.mere_id
            AND al.type = 'velage'
            AND DATE(al.date_creation) = CURDATE()
      );
END$$
DELIMITER ;
\end{sqlbox}

\begin{sqlbox}{evt\_rapport\_croissance (hebdomadaire)}
DELIMITER $$
CREATE EVENT evt_rapport_croissance
ON SCHEDULE EVERY 1 WEEK
STARTS CURRENT_TIMESTAMP
DO
BEGIN
    DECLARE v_nb_animaux INT;
    SELECT COUNT(*) INTO v_nb_animaux FROM animaux WHERE statut = 'actif';

    INSERT INTO alertes (animal_id, type, message, niveau)
    VALUES (NULL, 'autre',
        CONCAT('Rapport hebdo : ', v_nb_animaux,
               ' animaux actifs. Consultez le tableau de bord.'),
        'info');
END$$
DELIMITER ;
\end{sqlbox}

\begin{sqlbox}{evt\_alerte\_cout\_mensuel (mensuel)}
DELIMITER $$
CREATE EVENT evt_alerte_cout_mensuel
ON SCHEDULE EVERY 1 MONTH
STARTS '2026-04-01 00:00:00'
DO
BEGIN
    DECLARE v_total DECIMAL(12,2);

    SELECT SUM(quantite_kg * cout_unitaire_kg) INTO v_total
    FROM alimentation
    WHERE date_alimentation >= DATE_SUB(CURDATE(), INTERVAL 1 MONTH);

    INSERT INTO alertes (animal_id, type, message, niveau)
    VALUES (NULL, 'autre',
        CONCAT('Bilan mensuel alimentation : ',
               COALESCE(ROUND(v_total, 0), 0), ' FCFA.'),
        'info');
END$$
DELIMITER ;
\end{sqlbox}

% ────────────────────────────────────────────────────────────
\chapter{SYSTEM\_PROMPT complet}
\label{app:sysprompt}
% ────────────────────────────────────────────────────────────

Le SYSTEM\_PROMPT suivant est défini dans \texttt{app.py} via la fonction
\texttt{get\_system\_prompt()}. La variable \texttt{\{today\}} est injectée
dynamiquement à chaque appel.

\begin{sqlbox}{SYSTEM\_PROMPT — app.py (get\_system\_prompt)}
Tu es BoviBot, l'assistant IA expert d'un elevage bovin.
Tu aides l'eleveur a gerer son troupeau en traduisant ses demandes
en SQL ou en actions.
Nous sommes le {today}.

[DB_SCHEMA]
Tables MySQL :
- races(id, nom, origine, poids_adulte_moyen_kg, production_lait_litre_jour)
- animaux(id, numero_tag, nom, race_id, sexe[M/F], date_naissance,
  poids_actuel, statut[actif/vendu/mort/quarantaine], mere_id, pere_id)
- pesees(id, animal_id, poids_kg, date_pesee, agent)
- sante(id, animal_id, type, date_acte, veterinaire, prochain_rdv)
- reproduction(id, mere_id, pere_id, date_velage_prevue, statut)
- alimentation(id, animal_id, type_aliment, quantite_kg, cout_unitaire_kg)
- ventes(id, animal_id, acheteur, prix_fcfa, date_vente)
- alertes(id, animal_id NULL, type, message, niveau, traitee)
- historique_statut(id, animal_id, ancien_statut, nouveau_statut)
- production_lait(id, animal_id, date_traite, quantite_litre, periode)
- stocks(id, nom, categorie, quantite_disponible, unite, seuil_alerte)
Fonctions : fn_age_en_mois(id), fn_gmq(id),
            fn_cout_total_elevage(id), fn_rentabilite_estimee(id)
Procedures : sp_enregistrer_pesee, sp_declarer_vente,
             sp_rapport_nutritionnel

### BLOC 1 - FORMAT DE REPONSE OBLIGATOIRE
Reponds exclusivement en JSON pur, sans markdown.
- Consultation : {"type":"query","sql":"SELECT...","explication":"..."}
- Action : {"type":"action","action":"nom","params":{...},"confirmation":"..."}
- Info : {"type":"info","sql":null,"explication":"..."}

### BLOC 2 - REGLES SQL
- Filtrer statut='actif' par defaut sauf demande explicite.
- Utiliser fn_age_en_mois(a.id), fn_gmq(a.id),
  fn_cout_total_elevage(a.id), fn_rentabilite_estimee(a.id).
- Ne jamais generer DELETE ou DROP. Limiter a 100 resultats.

### BLOC 3 - PROCEDURES DISPONIBLES
1. sp_enregistrer_pesee(animal_id, poids_kg, date, agent)
   Mots-cles : "enregistre pesee", "pese", "nouveau poids"
2. sp_declarer_vente(animal_id, acheteur, telephone, prix_fcfa,
                     poids_vente_kg, date_vente)
   Mots-cles : "declare vente", "vends", "cede l'animal"
3. sp_rapport_nutritionnel(animal_id)
   Mots-cles : "rapport nutritionnel", "consommation", "ration"
IMPORTANT : date au format YYYY-MM-DD. Jamais CURDATE() dans les params.

### BLOC 4 - RESOLUTION DES IDENTIFIANTS
- Consultation : WHERE numero_tag = 'TAG-XXX'.
- Action : passe le numero_tag comme animal_id, le backend resoudra.
- INTERDIT : sous-requetes SQL comme valeur de parametre.
- Si animal inconnu : type 'info' pour demander clarification.
\end{sqlbox}

% ────────────────────────────────────────────────────────────
\chapter{Déclaration d'usage de l'IA}
\label{app:ia_declaration}
% ────────────────────────────────────────────────────────────

Conformément à la consigne §12 du Cahier des Charges, le groupe déclare l'usage
suivant des outils d'intelligence artificielle :

\begin{table}[H]
\caption{Déclaration d'usage des outils IA — \bovibot{}}
\label{tab:ia_declaration}
\small\bovirowcolors
\begin{tabularx}{\textwidth}{@{}L{2.8cm}L{4.5cm}X@{}}
  \toprule
  \rowcolor{bovibrown!15}
  \textbf{Outil} & \textbf{Tâches réalisées} & \textbf{Niveau d'intervention humaine} \\
  \midrule
  \textbf{Claude} (Anthropic) &
    Rédaction et structuration du rapport ; génération de squelettes PL/SQL ; justifications métier &
    Supervision complète — chaque section relue, corrigée et validée. Code testé en base avant intégration. \\
  \textbf{DeepSeek Chat} &
    Modèle LLM utilisé en production dans \bovibot{} pour le Text-to-SQL &
    Intégré via API — outputs validés par \texttt{validate\_sql()}. SYSTEM\_PROMPT rédigé par le groupe. \\
  \textbf{ChatGPT} (OpenAI) &
    Débogage de requêtes SQL complexes ; génération de données de test &
    Validation humaine systématique avant insertion en base. \\
  \textbf{GitHub Copilot} &
    Complétion de code Python (FastAPI) et JavaScript &
    Mode suggestion uniquement — chaque suggestion acceptée après lecture. \\
  \textbf{Stitch AI} &
    Conception et design des pages de l'interface frontend (maquettes, mise en page, composants visuels) &
    Les wireframes et choix de design ont été supervisés et validés par le groupe — Stitch AI utilisé comme outil de prototypage rapide. \\
  \textbf{Gemini} (Google) &
    Affinage et optimisation du code frontend (JavaScript, composants UI) &
    Chaque suggestion de code relue, comprise et adaptée par le groupe avant intégration. \\
  \bottomrule
\end{tabularx}
\end{table}

\begin{successbox}
  \textbf{Déclaration d'intégrité :} L'architecture du système, les choix de
  modélisation, la stratégie de prompt engineering et l'ensemble des décisions
  techniques ont été entièrement conçus par les membres du groupe. Les outils IA
  ont servi d'assistants à la rédaction et à l'implémentation, non de substituts
  à la réflexion conceptuelle. Plus précisément, le raisonnement, les décisions
  de conception, les workflows, le travail de modélisation des données et
  l'ensemble de la réflexion intellectuelle ont été assurés par les membres du
  groupe. Les outils d'IA ont été mobilisés uniquement parce qu'ils sont plus
  rapides que nous dans l'écriture de code et de texte — ils n'ont en aucun cas
  remplacé la pensée, la stratégie, ni les choix d'architecture qui restent
  l'œuvre exclusive de l'équipe.
\end{successbox}

\end{document}
